% paper/equivalence.tex — Ex ante equivalence of the labour-closure endpoints
% under demand-only shocks (ALPHA ≡ BETA, GAMMA ≡ DELTA).
%
% Repo: BFRep/(3)BeyondHulten, branch `reorg`, base commit c222f7b (2026-09-18).
% Cited evidence: the matrix_5x3-v5-* runs (ADR-0019 generation, 15/15 executed
% 2026-09-18) and paper/tables/matrix_5x3_v5_flows.md; the cbase2 analytic
% equivalence test recorded in docs/DOCS_ASSESSMENT.md §4.2; the supply-side
% pilot numbers from docs/ETAs.md.
% Companion: docs/ETAs.md (the supply-side arm and the demand-sensitive-price
% candidates); paper/closures_and_demand_shocks.md (the CGE-vs-IO framing);
% ADR-0020 (proposed) for the eta = 0 external-position finding reflected in
% the BF remarks.
%
% No LaTeX in the analysis container: compile on the Mac (pdflatex; the
% bibliography below is a plain thebibliography, no .bib dependency).

\documentclass[11pt]{article}

\usepackage[english]{babel}
\usepackage[a4paper,top=2.2cm,bottom=2.2cm,left=3cm,right=3cm]{geometry}
\usepackage{amsmath,amssymb,amsthm}
\usepackage{booktabs}
\usepackage[colorlinks=true,allcolors=blue]{hyperref}

\setlength\parindent{0pt}
\setlength{\parskip}{0.45em}

\newtheorem{lemma}{Lemma}
\newtheorem{proposition}{Proposition}
\newtheorem{corollary}{Corollary}
\newtheorem{remark}{Remark}

\numberwithin{equation}{section}

\newcommand{\Lbar}{\bar{L}}
\newcommand{\one}{\mathbf{1}}

\title{Ex Ante Equivalence of the Labour-Closure Endpoints\\
under Demand-Only Shocks\\[0.4em]
\large ALPHA $\equiv$ BETA and GAMMA $\equiv$ DELTA in the \texttt{BeyondHulten} kernel}
\author{Hermes Agent (Lt.\ Cmdr.\ Data), for Prof.\ Dr.\ J.\ Kapeller}
\date{September 18, 2026}

\begin{document}
\maketitle

\begin{abstract}
The executed $5\times3$ closure matrix reports two exact coincidences: the
BETA row reproduces ALPHA bit-for-bit, and the DELTA row reproduces GAMMA
cell-for-cell, in every financing column. This note shows that both
coincidences are derivable \emph{ex ante} from the structure of the kernel,
states the conditions under which they hold exactly, and classifies what they
mean for the manuscript: under a demand-only programme the price block is
demand-invariant (a one-factor Samuelson nonsubstitution property), so the
equilibrium real wage sits at its technology-determined anchor and the whole
quantity system reduces to one linear Leontief multiplier system in which none
of the elasticity parameters $(\theta,\epsilon,\sigma,\eta_s)$ appears. The
four closure rows differ only in which scalar absorbs the shock: the external
transfer $F$ (mobile rows, employment pinned), employment $L$ (fixed-wage
rows, $F \equiv 0$), or neither (BF, no induced consumption round). The final
section derives what would break the equivalences --- the analytic basis of
the demand-sensitive-price program of \texttt{docs/ETAs.md} --- and proves a
no-go corollary: nominal wage rules cannot create price variety in this
framework.
\end{abstract}

\section{Purpose and summary}

Reviewer 2's critique turns on the exact conditions under which the closure
rows of the evaluation matrix are or are not equivalent
(\texttt{docs/DOCS_ASSESSMENT.md}, section 4.2). The executed
\texttt{matrix\_5x3\_v5} generation (ADR-0019; 15/15 cells, run ids
\texttt{matrix\_5x3-v5-*}) measures two coincidences:

\begin{itemize}
\item \textbf{ALPHA $\equiv$ BETA.} The elastic-supply row
      ($\eta_s = 0.5$) reports every metric of the full-employment row
      identically: real GDP, the welfare index, employment, the external
      account --- to all eight printed decimals, with the identity wedge at
      $\sim 10^{-16}$. A direct comparison of stored solutions measured
      $\max|\Delta y| = 6.2\times 10^{-17}$ (session of 2026-09-17).
\item \textbf{GAMMA $\equiv$ DELTA.} The Leontief-corner row
      ($( \theta_\delta, \epsilon_\delta, \sigma_\delta ) = (10^{-4},
      10^{-4}, 10^{-4})$) reports every metric of the fixed-wage row at the
      baseline CES elasticities $(\theta,\epsilon,\sigma) = (0.5,0.5,0.9)$
      identically, in all three financing columns; the flow-table wedges
      agree at the residual level ($\sim 10^{-12}$). The analytic
      \texttt{leontief\_multiplier} test reproduces the nonlinear DELTA
      solution with relative error $0.0$ (recorded in
      \texttt{docs/DOCS\_ASSESSMENT.md} section 4.2).
\end{itemize}

Sections~\ref{sec:kernel}--\ref{sec:reduction} set up the common kernel and
derive the demand-only reduction. Section~\ref{sec:equiv} states and proves
the two equivalences and classifies them. Section~\ref{sec:gauge} isolates
what \emph{does} differ across the rows (the scalar closure regimes).
Section~\ref{sec:scope} derives the scope conditions: which model changes
break the equivalences --- the menu behind the demand-sensitive-price program
--- and which cannot.

Throughout, ``demand-only shock'' means the matrix design: a programme bundle
$g = m\,G_0\,\psi$ entering final demand under one of the three financing
closures, with the supply shock at $A \equiv 1$ (unit productivity).

\section{The common kernel}\label{sec:kernel}

Let $N$ be the number of sectors. Prices $p_i$, wage $w$, gross output
$y_i$, employment $L_i$. Data-calibrated objects: factor shares $\alpha_i$
(labour), domestic bill coefficients $a_u = A^{bill}_u/\lambda_u$ with
supplier-share matrix $\Omega$ ($\Omega_{iu}$ = sector $i$'s share of input
$u$; rows sum to one), CPI weights $\omega_i$ ($\sum_i \omega_i = 1$),
sector import margins $m_i$, saving rate $s$, exogenous government demand
$g^G$, investment $inv$, exports $X$ (no margin), the programme bundle $g$,
and the labour bar $\Lbar = \sum_i \bar{L}_i = 1$ (baseline income units,
GDP $=1$). Elasticities: intermediate substitution $\theta$,
labour--intermediate substitution $\epsilon$, consumption substitution
$\sigma$, allocation $\eta \in \{0,1\}$ (ADR-0010), supply elasticity
$\eta_s \geq 0$ (BETA only).

\paragraph{Price block.} The CES intermediate price index and unit cost are
\begin{align}
P_i(p) &= \Big( \textstyle\sum_u \Omega_{iu}\, p_u^{\,1-\theta}
   \Big)^{\frac{1}{1-\theta}}, \label{eq:ip}\\
c_i(p,w;A) &= \Big( A_i^{\,\epsilon-1}\big( \alpha_i\, w^{1-\epsilon}
   + (1-\alpha_i)\, P_i^{\,1-\epsilon} \big) \Big)^{\frac{1}{1-\epsilon}},
   \label{eq:cost}
\end{align}
with the Cobb--Douglas limit $c_i = w^{\alpha_i} P_i^{\,1-\alpha_i}/A_i$ at
$\epsilon = 1$. Zero profit requires
\begin{equation}
p_i = c_i(p,w;A) \qquad \forall\, i. \label{eq:zp}
\end{equation}
At the calibration point $(p,w,A) = (\one,1,\one)$ the unit cost equals one
exactly, because $\sum_u \Omega_{iu} = 1$ makes $P_i = 1$; this is the
calibration identity used throughout.

\paragraph{Quantity block.} With full cost-minimizing allocation
($\eta = 1$) the marginal-product condition inverts to
\begin{equation}
L_i = \big( p_i\, A_i^{(\epsilon-1)/\epsilon}\, \alpha_i^{1/\epsilon}\,
      y_i^{1/\epsilon} / w \big)^{\epsilon}
    = p_i^{\epsilon}\, A_i^{\epsilon-1}\, \alpha_i\, y_i / w^{\epsilon},
\label{eq:labdemand}
\end{equation}
At $\eta = 0$ (the BF endpoint) the allocation is the frozen baseline,
$L_i = \bar{L}_i$. Intermediate demand is charged with the \emph{domestic}
bill (ADR-0012):
\begin{equation}
int_i = p_i^{-\theta} \sum_u \Omega_{ui}\,
        p_u^{\epsilon} A_u^{\epsilon-1} P_u^{\theta-\epsilon}\, a_u\, y_u.
\label{eq:int}
\end{equation}
Household demand is CES with CPI deflation and the financing weights
$ds$ ($ds = 1$ except under F1, where $ds = d$, the preference tilt):
\begin{equation}
c_i = (1-s)(1-m_i)\,\omega_i\, ds_i\; E\; p_i^{-\sigma} \Big/
      \textstyle\sum_j \omega_j ds_j\, p_j^{1-\sigma},
\qquad
E = (1-\tau)\, w \textstyle\sum_i L_i + F,
\label{eq:cons}
\end{equation}
with the lump-sum tax rate $\tau = p\!\cdot\! g^G / (w \sum_i L_i)$ under
F1/F3 and $\tau = p\!\cdot\!(g^G + g)/(w \sum_i L_i)$ under F2, and $F$ the
net external transfer entering expenditure after tax (ADR-0019). Saving
$sE$ leaks; the import content is charged margin by margin. Goods-market
clearing reads
\begin{equation}
y_i = int_i + c_i + (1-m_i)\big( g^G_i + inv_i + g_i \big) + X_i,
\label{eq:clearing}
\end{equation}
where $g_i$ appears only under the additive financings F2/F3
($g_i = 0$ under F1; legacy manna is zero in the matrix).

\paragraph{The four regimes.} The mobile systems solve for
$X = (p, y, w, F) \in \mathbb{R}^{2N+2}$ with the numeraire
$\Pi(p) = 1$, where
\begin{equation}
\Pi(p) = \Big( \textstyle\sum_i \omega_i\, p_i^{1-\sigma}
   \Big)^{\frac{1}{1-\sigma}},
\label{eq:cpi}
\end{equation}
and close with the labour equation
\begin{align}
\text{ALPHA:}\quad & \textstyle\sum_i L_i(p,y,w) = \Lbar,
\label{eq:alpha}\\
\text{BETA:}\quad & \textstyle\sum_i L_i(p,y,w)
   = \Lbar \Big[ \frac{w/\Pi(p)}{w_0/\Pi_0} \Big]^{\eta_s},
\qquad w_0/\Pi_0 = 1.
\label{eq:beta}
\end{align}
The fixed-wage systems solve for $X = (p, y) \in \mathbb{R}^{2N}$ with the
wage pinned $w \equiv 1$ and $F \equiv 0$; employment is a post-solve
outcome. GAMMA runs them at the baseline elasticities, DELTA at the Leontief
corner $(\theta_\delta,\epsilon_\delta,\sigma_\delta) = (10^{-4},10^{-4},
10^{-4})$. BF is the mobile system at $\eta = 0$ with the $F = 0$ pin
(decision D2 of ADR-0019). All rows in the matrix use $\eta = 1$ except BF.

\section{The demand-only reduction}\label{sec:reduction}

\begin{lemma}[Price block demand-invariance]\label{lem:price}
The zero-profit system \eqref{eq:zp} is homogeneous of degree one in
$(p,w)$ jointly and contains no demand variable. Consequently, for any
technology $A$, every solution satisfies $p = w\,\pi(A)$, where
$\pi(A) \in \mathbb{R}^N_{++}$ solves $c_i(\pi,1;A) = \pi_i$ for all $i$;
relative prices and the ratio $p/w$ are functions of $A$ alone. Under the
CPI numeraire, $w = 1/\Pi(\pi(A))$ and $p = \pi(A)/\Pi(\pi(A))$; under the
fixed wage, $w = 1$ and $p = \pi(A)$. At $A \equiv \one$ the calibration
identity gives $\pi = \one$, hence in both gauges
\begin{equation}
p = \one, \qquad w = 1, \qquad \Pi = 1,
\end{equation}
independently of the demand vector, of the financing closure, and of the
shock size.
\end{lemma}

\begin{proof}
Homogeneity: $c_i(tp, tw; A) = t\, c_i(p,w;A)$ for $t>0$, by inspection of
\eqref{eq:cost}--\eqref{eq:ip}. Write $q = p/w$. Then \eqref{eq:zp} is
equivalent to $q_i = c_i(q,1;A)$ for all $i$, i.e.\ $q = \pi(A)$, and the
solution set of \eqref{eq:zp} is exactly the ray $\{(w\,\pi(A), w)\}$. The
CPI \eqref{eq:cpi} is homogeneous of degree one in $p$, so
$\Pi(w\,\pi) = w\,\Pi(\pi)$; the numeraire $\Pi = 1$ or the pin $w = 1$
selects the scale. At $A \equiv \one$, $\pi = \one$ solves
\eqref{eq:zp} at $w = 1$ because $c_i(\one,1;\one) = 1$ (the calibration
identity), and $\Pi(\one) = 1$ because $\sum_i \omega_i = 1$.
\end{proof}

\begin{corollary}[Real-wage invariance]\label{cor:realwage}
At every solution of the shared block,
\begin{equation}
\frac{w}{\Pi(p)} = \frac{1}{\Pi(\pi(A))},
\end{equation}
a function of technology alone. At $A \equiv \one$ this equals $1 =
w_0/\Pi_0$: the equilibrium real wage sits at its anchor for every demand
vector, every financing closure and every shock magnitude. (Measured in the
matrix: the GDP deflator reads $1.000000$ in all 15 cells, and a demand
shock of ten times the programme size moves $\max|p-\one|$ by
$2.7\times10^{-15}$; \texttt{docs/ETAs.md}.)
\end{corollary}

\begin{lemma}[Elasticity-free affine quantity block]\label{lem:affine}
Fix $\eta = 1$ and $A \equiv \one$, and evaluate the quantity block at the
demand-only prices $(p,w) = (\one,1)$. Then for \emph{every} value of
$(\theta,\epsilon,\sigma)$:
\begin{align}
int &= \Omega^{\!\top}\mathrm{diag}(a)\, y, &
L &= \mathrm{diag}(\alpha)\, y, &
\sum_i L_i &= \alpha^{\!\top} y,
\label{eq:leontiefblocks}\\
c &= (1-s)\,\mathrm{diag}(1-m)\,\tilde\omega(ds)\; E, &
E &= \alpha^{\!\top} y - \textstyle\sum g^G + F
   \;\; (\text{F1/F3}), \nonumber
\end{align}
where $\tilde\omega(ds) = (\omega \odot ds)\,/\, \one^{\!\top}(\omega \odot
ds)$ and $E = \alpha^{\!\top} y - \sum g^G - \sum g$ under F2. Hence the
clearing system \eqref{eq:clearing} is exactly affine,
\begin{equation}
y = G\, y + b(F),
\qquad
G = \Omega^{\!\top}\mathrm{diag}(a)
  + (1-s)\,\mathrm{diag}(1-m)\,\tilde\omega(ds)\,\alpha^{\!\top},
\label{eq:affine}
\end{equation}
\begin{equation}
b(F) = (1-m)\odot\big( g^G + inv + g \big) + X
     + (1-s)\,\mathrm{diag}(1-m)\,\tilde\omega(ds)
       \big( F - \textstyle\sum g^G - \textstyle\sum g \big),
\label{eq:bvec}
\end{equation}
with $g$ present only under F2/F3 and the last bracket reduced accordingly.
No elasticity parameter enters $G$ or $b$. At $\eta = 0$ (BF),
$\sum_i L_i = \sum_i \bar{L}_i = \Lbar$ is a constant, $E$ does not depend
on $y$, and the household round drops out of $G$:
$G_{BF} = \Omega^{\!\top}\mathrm{diag}(a)$.
\end{lemma}

\begin{proof}
At $p = w = A = \one$: $P_u = 1$ in \eqref{eq:int}, so all elasticity
factors $p^{\epsilon}$, $P^{\theta-\epsilon}$, $p^{-\theta}$ equal $1$ and
\eqref{eq:int} reduces to the fixed-coefficient round
$\sum_u \Omega_{ui} a_u y_u$; \eqref{eq:labdemand} reduces to
$L_i = \alpha_i y_i$ since every exponent collapses at unit prices; in
\eqref{eq:cons} the terms $p_i^{-\sigma}$ and the normalizer reduce to
$1$ and $\sum_j \omega_j ds_j$, giving the fixed spending weights
$\tilde\omega(ds)$; and the tax rate $\tau = \sum g^G / \alpha^{\!\top}y$
turns $E$ into the affine form shown. Collecting terms in
\eqref{eq:clearing} gives \eqref{eq:affine}--\eqref{eq:bvec}: the
intermediate round contributes the columns
$\sum_i \Omega_{iu} a_u = a_u$ (row sums of $\Omega$ equal $1$) and the
household round contributes $(1-s)\sum_i (1-m_i)\tilde\omega_i(ds)\,
\alpha_u$ per unit of $y_u$.
\end{proof}

\begin{remark}[Nonsubstitution]\label{rem:samuelson}
Lemma~\ref{lem:price} and Lemma~\ref{lem:affine} are the production side of
the Samuelson nonsubstitution theorem: with one primary factor, constant
returns and no joint production, the cost-minimizing technique --- and hence
relative prices --- is independent of the demand vector, and final demand
allocates quantities only (Samuelson 1951). The kernel's nested-CES
structure is a smooth parametrisation of that corner; the demand-only matrix
therefore inherits the theorem's rigidity, and this is what
\texttt{docs/DOCS\_ASSESSMENT.md} section 4.2 records as the ``model-class
property'' behind the coincidences. The IO endpoint makes the rigidity
literal --- the Robinson (2006) fixed-price multiplier --- but, as
Proposition~\ref{prop:gd} shows, it is not an approximation: \emph{every}
elasticity triple in the fixed-wage gauge lives on the same corner.
\end{remark}

\section{The two equivalences}\label{sec:equiv}

\begin{proposition}[ALPHA $\equiv$ BETA]\label{prop:ab}
Let $A \equiv \one$ and $\eta = 1$. The ALPHA system
\eqref{eq:zp},\eqref{eq:clearing},\eqref{eq:alpha},$\Pi = 1$ and the BETA
system \eqref{eq:zp},\eqref{eq:clearing},\eqref{eq:beta},$\Pi = 1$ have
identical solution sets, for every $\eta_s \geq 0$ and every financing
closure.
\end{proposition}

\begin{proof}
The two systems differ in exactly one equation. By
Lemma~\ref{lem:price} and Corollary~\ref{cor:realwage}, every solution of
the shared block satisfies $w/\Pi = 1 = w_0/\Pi_0$, so BETA's labour
equation \eqref{eq:beta} reduces to $\sum_i L_i = \Lbar$, which is
\eqref{eq:alpha}. Conversely, any solution of the ALPHA system satisfies
BETA's equation by the same substitution. The systems are therefore
equivalent, equation by equation.
\end{proof}

\begin{proposition}[GAMMA $\equiv$ DELTA]\label{prop:gd}
Let $A \equiv \one$, $\eta = 1$, $F \equiv 0$, $w \equiv 1$. Then
$(p,y) = (\one,\, y^{\ast})$ with
$y^{\ast} = (I - G)^{-1} b(0)$ solves the fixed-wage system at
\emph{every} elasticity triple $(\theta,\epsilon,\sigma)$; the employment
outcome is $L^{\ast} = \alpha^{\!\top} y^{\ast}$, the consumption block is
$c^{\ast} = (1-s)\,\mathrm{diag}(1-m)\,\tilde\omega(ds)\,
(L^{\ast} - \sum g^G - \sum g\,\mathbf{1}_{F2})$, where
$\mathbf{1}_{F2}$ selects the tax-financed column, and the block
decomposition \eqref{eq:leontiefblocks} holds verbatim. In particular GAMMA (baseline
elasticities) and DELTA (the Leontief corner) share this solution, and the
analytic system \eqref{eq:affine}--\eqref{eq:bvec} ---
\texttt{leontief\_multiplier} in \texttt{src/closures/labor/labor.jl} ---
computes it exactly.
\end{proposition}

\begin{proof}
Existence and form: at $p = \one$ the price block holds by the calibration
identity (Lemma~\ref{lem:price}), and by Lemma~\ref{lem:affine} the
clearing block at $p = \one$ is the affine system, of which
$y^{\ast} = (I-G)^{-1}b(0)$ is the solution. Uniqueness on the admitted
branch: the scale-determinacy guard (ADR-0017) verifies at the reference
prices that $(I - G)$ is nonsingular --- contraction if
$\max_u \mathrm{colsum}_u(G) < 1$, otherwise a direct rank test --- and
that the candidate $y^{\ast}$ is strictly positive; within the admitted
branch the solution is unique, and the multi-init admissibility checks
recorded for the F1 column (three distinct starting points converging to
the same root) support root selection. Both the GAMMA and the DELTA solve
are subject to the same guard, and both converged to this root (residuals
$6.7\times10^{-14}$ to $4.4\times10^{-16}$ across the six fixed-wage cells,
\texttt{runs/matrix\_5x3-v5-\{GAMMA,DELTA\}-*/manifest.toml}). The
elasticity-independence is exact in the equations, not an $\epsilon$-limit
artefact: the DELTA solve at $\epsilon_\delta = 10^{-4}$ reproduces the
analytic $y^{\ast}$ with relative error $0.0$ (F2/F3 verification, recorded
in \texttt{docs/DOCS\_ASSESSMENT.md} section 4.2), and the v5 flow metrics
of the two rows agree to all printed decimals with residual-level wedge
differences $\sim 10^{-12}$.
\end{proof}

\begin{corollary}[Classification]\label{cor:class}
For the manuscript's framing of the matrix:

\begin{center}
\begin{tabular}{lll}
\toprule
Coincidence & Status & What is unidentified \\
\midrule
GAMMA $\equiv$ DELTA & definitional in form, & $(\theta,\epsilon,\sigma)$ \\
 & substantive in mechanism & in the fixed-wage gauge \\
 & (Lemma~\ref{lem:affine}) & \\
ALPHA $\equiv$ BETA & derivable & $\eta_s$ in the flexible-wage gauge \\
 & (Cor.~\ref{cor:realwage}) & \\
F2 $\equiv$ F3 (mobile rows) & derivable & the financing label; \\
 & (Section~\ref{sec:gauge}) & $F_{F3} = F_{F2} - B_{gov}$ \\
BF vs ALPHA & pin artifact, not an allocation & the $\eta=0$ external \\
 & effect (ADR-0020 audit: the labour- & position is not \\
 & clearing closure reproduces ALPHA & identified; never \\
 & to $2.8\times10^{-17}$) & print it \\
\bottomrule
\end{tabular}
\end{center}

None of these is a numerical accident; all four hold at machine precision
in the executed generation, and none is a mis-specification. They are
properties of the \emph{experiment}: a demand-only programme in a one-factor
CRS economy cannot identify any elasticity of the kernel (the mechanism is
Remark~\ref{rem:samuelson}).
\end{corollary}

\begin{remark}[Off the demand-only slice]\label{rem:supply}
The equivalences fail the moment technology moves. At $A \neq \one$,
Corollary~\ref{cor:realwage} gives $w/\Pi = 1/\Pi(\pi(A)) \neq 1$ in
general, and BETA's labour equation reduces to
$\alpha^{\!\top}y = \Lbar\,[1/\Pi(\pi(A))]^{\eta_s}$: employment is pinned
to a \emph{technology-determined constant} that now depends on $\eta_s$.
The supply-side pilot (+20\% in sector 1) measures exactly this separation
--- the failure of Proposition~\ref{prop:ab}'s hypothesis $A \equiv \one$:
$w = 1.004903$, with $L = 1.0024$ at $\eta_s = 0.5$ versus $L = 1.0098$ at
$\eta_s = 2$ against ALPHA's $L = 1.0000$ (\texttt{docs/ETAs.md}). The
identification of $\eta_s$ therefore requires a supply-side scenario ---
the premise of the supply arm.
\end{remark}

\section{What actually differs across the rows}\label{sec:gauge}

Lemma~\ref{lem:affine} shows that the demand-only matrix is a partition of
\emph{scalar closure regimes} over one and the same linear system. Three
regimes occur:

\begin{center}
\begin{tabular}{lccc}
\toprule
Rows & $F$ & Employment & Household round in $G$ \\
\midrule
ALPHA, BETA ($\eta = 1$) & free (instrument) & pinned, $L = \Lbar$ & active \\
GAMMA, DELTA ($\eta = 1$) & $F \equiv 0$ & endogenous, $L = \alpha^{\!\top}y$ & active \\
BF ($\eta = 0$) & $F \equiv 0$ & pinned via allocation & absent \\
\bottomrule
\end{tabular}
\end{center}

\paragraph{The $F$-instrument.} In the mobile regime the labour pin selects
the transfer. Writing $y(F) = (I-G)^{-1}(b_0 + f\,F)$ with
$f = (1-s)\,\mathrm{diag}(1-m)\,\tilde\omega(ds)$, the pin
$\alpha^{\!\top}y = \Lbar$ gives
\begin{equation}
F^{\ast} = \frac{\Lbar - \alpha^{\!\top}(I-G)^{-1} b_0}
               {\alpha^{\!\top}(I-G)^{-1} f}.
\label{eq:finstrument}
\end{equation}
This instrument is why the mobile F2 and F3 columns are real-neutral: F2
charges the programme through the lump-sum tax ($E = \alpha^{\!\top}y -
\sum g^G - \sum g$), F3 through the external account ($E = \alpha^{\!\top}y
- \sum g^G + F$), and the two $b$-vectors differ by exactly
$f\,\sum g$. The pin \eqref{eq:finstrument} therefore yields
$F^{\ast}_{F3} = F^{\ast}_{F2} - \sum g$, identical real allocations, and
an identical net external position $F + B_{gov}$. Measured: $F_{F3} - F_{F2} = -0.01330991 = -G_0$ to eight
decimals in every mobile row (Table~\ref{tab:v5}).

\paragraph{Who pays.} The economically meaningful difference between the
two paradigms is which margin absorbs the programme. In the mobile rows the
labour pin forces the household to pay (consumption $-1.82\%$ under F2/F3;
the $F$ instrument crowds it out exactly), while real GDP is flat by
construction ($w\,\Lbar$ deflated, ADR-0018). In the fixed-wage rows
employment is the margin: F3 expands employment by $1.78\%$ and consumption
by $2.26\%$ (the Keynesian case), F2 leaves GDP nearly neutral
($+0.13\%$) with a welfare loss from the tax, and F1 shifts composition
slightly downward ($-0.10\%$ employment). BF is the degenerate third
regime: with income constant and the household round absent, F3 leaves the
welfare index at exactly $0.000000$. Two caveats from the ADR-0020 audit
(proposed; \texttt{docs/log/2026-09.md}): the recorded BF/ALPHA differences
are carried by the pin, not by the frozen allocation --- replacing the pin
by the labour equation reproduces the ALPHA cells exactly
($2.8\times10^{-17}$) and extends financing neutrality to $\eta = 0$ --- and
the $\eta = 0$ external position is not identified (every pin value $F = c$
is an exact root that moves the booked position one-for-one), so no BF
external position should be printed.

\paragraph{Admissibility.} All fifteen cells execute under the v5
calibration. In particular the F1 column is admissible in the fixed-wage
rows: the endogenous tax rate $\tau = \sum g^G/\alpha^{\!\top}y$ supplies
the marginal leakage that keeps $\max \mathrm{colsum}(G) < 1$, as verified
by the ADR-0017 criterion (the retired heuristic guard had blocked these
cells in the v1 generation).

\section{Scope conditions: what breaks the equivalences}\label{sec:scope}

Both propositions rest on Lemma~\ref{lem:price}: the price block contains
no demand variable. Any model change that lets demand into that block
restores identification and creates variety in the published matrix. The
doors are:

\begin{enumerate}
\item \textbf{Labour door --- sector-specific wages.} With $N$ sectoral
      labour markets $w_i = \phi_i(L_i)$ instead of one common $w$, the
      homogeneity argument of Lemma~\ref{lem:price} collapses: $L_i =
      \alpha_i y_i$ depends on demand, so $w_i$ and through zero profit
      $p_i$ become demand-sensitive. Zero profit is preserved (no rents),
      the new parameters are sectoral supply elasticities of the same
      family as $\eta_s$, and the variant repairs the BF row's
      common-wage defect (ADR-0017's open gate). Ranked candidate 1 in
      \texttt{docs/ETAs.md}.
\item \textbf{Markup door --- demand-sensitive markups.} A pricing rule
      $p_i = \mu_i(y_i)\, c_i(p,w)$, e.g.\ $\mu_i = 1 + \kappa\,(y_i/y_{i0}
      - 1)$, breaks Lemma~\ref{lem:price} directly. The real cost is not
      code but accounting: markup income is profit, so a
      profit-income-and-spending closure is required or the external-account
      identity behind the canary breaks. Nested at $\kappa = 0$. Candidate 2
      in \texttt{docs/ETAs.md}.
\item \textbf{Capacity door --- a fixed factor, or its reduced form.} A
      sector-fixed factor with a rental schedule makes marginal cost rise
      in $y_i$ (rents require a rental-income closure --- the largest
      surgery, candidate 3 in \texttt{docs/ETAs.md}). The reduced-form
      variant avoids all income-closure problems: an endogenous utilization
      externality $A^{eff}_i = A_i\,(y_i/\lambda_i)^{-\delta}$ raises unit
      cost with demand while keeping zero profit (the extra cost is paid to
      no one), so a one-line change in the cost hook \eqref{eq:cost}
      suffices. With the opposite sign ($\delta < 0$) it is the
      Kaldor--Verdoorn case: increasing returns make prices \emph{fall}
      with demand --- heterodox-friendly variety in the other direction.
      (New candidate; see Section~\ref{sec:conclusion}.)
\item \textbf{External door --- endogenous import prices.} An
      upward-sloping foreign supply curve
      $p^M_i = p^M_{i0}\,(M_i/M_{i0})^{\zeta}$, $\zeta \geq 0$, enters the
      intermediate index \eqref{eq:ip} and the CPI through the margins, so
      $\pi$ becomes a function of import volumes and hence of demand. The
      surgery is moderate and it exploits precisely the external-account
      machinery the paper has just built (ADR-0019): the programme's import
      content would bid up import prices --- a structuralist
      absorption/term-of-trade channel that fits the open-economy framing.
      (New candidate.)
\item \textbf{Technology door --- supply shocks.} $A \neq 1$ breaks
      Corollary~\ref{cor:realwage} and separates BETA from ALPHA
      (Remark~\ref{rem:supply}); this is the supply arm already designed in
      \texttt{docs/ETAs.md}, not a repair of the demand-only matrix.
\end{enumerate}

\begin{corollary}[No-go for nominal wage rules]\label{cor:nogo}
Within the kernel's one-factor CRS price block, no rule that writes the
\emph{nominal} wage as a function of endogenous aggregates --- an aggregate
wage curve $w = \bar{w}\,(\alpha^{\!\top}y/\Lbar)^{\varphi}$, nominal
indexation, or a re-anchored numeraire --- can create demand-sensitive
\emph{real} outcomes. By Lemma~\ref{lem:price}, $p = w\,\pi(A)$ and
$\Pi(p) = w\,\Pi(\pi)$ at every solution, so the real wage
$w/\Pi = 1/\Pi(\pi(A))$ is invariant to whatever sets $w$; the rule merely
re-labels the nominal scale. With the CPI numeraire or the $w = 1$ pin such
a rule is either vacuous or over-determines the system (in the fixed gauge
it forces $\alpha^{\!\top}y = \Lbar$, i.e.\ silently reverts GAMMA to
ALPHA). A rule on the \emph{real} wage, $w/\Pi = \phi(L)$, is BETA's
equation \eqref{eq:beta} rearranged and is already in the model --- where
it is pinned by technology under demand-only shocks. The numeraire lever is
likewise dead (DE-0011).
\end{corollary}

Quantity-side levers (endogenous saving rates, endogenous export demand)
change $G$ or $b$ but not the price block; they add variety between
financing columns, not between closures, and leave Lemma~\ref{lem:price}
intact.

\section{Numerical evidence}\label{sec:conclusion}

Table~\ref{tab:v5} extracts the four rows of the executed
\texttt{matrix\_5x3\_v5} generation (income-side real GDP and the
household-consumption welfare index, ADR-0018; employment; external
account). The paired rows are identical to all printed decimals in every
reported metric, including the full flow decomposition (saving, tax,
investment plus exports, imports, taxes on intermediates; see
\texttt{paper/tables/matrix\_5x3\_v5\_flows.md}). The GDP deflator reads
$1.000000$ in all fifteen cells, the numerical face of
Lemma~\ref{lem:price}.

\begin{table}[htbp]
\centering
\caption{The equivalence rows in the \texttt{matrix\_5x3\_v5} generation
(run ids \texttt{matrix\_5x3-v5-*}; deviations from the no-programme
baseline).}
\label{tab:v5}
\begin{tabular}{llrrr}
\toprule
Row & Financing & Real GDP & Consumption (welfare) & Employment $L$ \\
\midrule
ALPHA & F1 & $+0.000000$ & $+0.001432$ & $1.000000$ \\
BETA  & F1 & $+0.000000$ & $+0.001432$ & $1.000000$ \\
ALPHA & F2 & $+0.000000$ & $-0.018226$ & $1.000000$ \\
BETA  & F2 & $+0.000000$ & $-0.018226$ & $1.000000$ \\
ALPHA & F3 & $+0.000000$ & $-0.018226$ & $1.000000$ \\
BETA  & F3 & $+0.000000$ & $-0.018226$ & $1.000000$ \\
\midrule
GAMMA & F1 & $-0.000973$ & $-0.000806$ & $0.999027$ \\
DELTA & F1 & $-0.000973$ & $-0.000806$ & $0.999027$ \\
GAMMA & F2 & $+0.001260$ & $-0.015333$ & $1.001260$ \\
DELTA & F2 & $+0.001260$ & $-0.015333$ & $1.001260$ \\
GAMMA & F3 & $+0.017796$ & $+0.022644$ & $1.017796$ \\
DELTA & F3 & $+0.017796$ & $+0.022644$ & $1.017796$ \\
\bottomrule
\end{tabular}
\end{table}

\paragraph{Implications for the manuscript.} (i) The two-paradigm reading
is safe to state as a theorem-backed claim: the flexible-wage rows form one
equivalence class (composition at pinned employment), the fixed-wage rows
another (employment at pinned real wage), and the difference between the
classes is the $F$-instrument --- a derivation, not an observation.
(ii) The null results for $\eta_s$ and for $(\theta,\epsilon,\sigma)$ under
demand-only shocks should be reported as identification statements
(Corollary~\ref{cor:class}), which is precisely what reviewer 2's critique
needs. (iii) The variety program for the demand-only matrix must route
through one of the five doors of Section~\ref{sec:scope}; within it, the
capacity door's utilization/Verdoorn variant and the external door are the
cheap additions beyond the three candidates already ranked in
\texttt{docs/ETAs.md}, while Corollary~\ref{cor:nogo} rules out the entire
nominal-wage-rule family ex ante.

\begin{thebibliography}{9}

\bibitem{samuelson1951}
Samuelson, P.\,A.\ (1951). Abstract of a theorem concerning substitutability
in open Leontief models. In T.\,C.\ Koopmans (ed.),
\emph{Activity Analysis of Production and Allocation}. Wiley, New York.

\bibitem{robinson2006}
Robinson, S.\ (2006). Macro models and multipliers: Leontief, Stone, Keynes,
and CGE models. IFPRI, Washington, DC.

\bibitem{baqaee2019}
Baqaee, D.\,R., and E.\ Farhi (2019). Beyond Hulten's theorem.
\emph{Econometrica} 87(3).

\bibitem{lavoie2014}
Lavoie, M.\ (2014). \emph{Post-Keynesian Economics: New Foundations}.
Edward Elgar, Cheltenham.

\end{thebibliography}

\end{document}
